\cleardoublepage
\thispagestyle{empty}
\begin{flushleft}
	\begin{tabular}{@{}l l@{}}
		\textbf{Titel:} &
		\parbox[t]{0.65\textwidth}{\textit{Wege der Analysis}}\\[0.4em]
		
		\textbf{Untertitel:} &
		\parbox[t]{0.65\textwidth}{Von der Bewegung eines Punktes zu Green, Stokes und Gauß}\\[1.2em]
		
		\textbf{Autor:} & Dipl.-Ing.\,(FH) Christian Weilharter\\[0.5em]
		\textbf{\copyright\ 2026} & Christian Weilharter, Traunstein\\[0.5em]
		\textbf{ISBN:} & xxxxxxxxxxxxxxx\\
		\textbf{Satz:} & \LaTeX\\
		\textbf{Druck:} & [Print-on-Demand-Dienst]\\
		\textbf{Kontakt:} & christian@weilharter.de\\
		\textbf{Web:} & www.mathandphysics.de\\
	\end{tabular}
\end{flushleft}


\vspace{2em}
\noindent
Alle Rechte vorbehalten. Kein Teil dieses Buches darf ohne schriftliche Genehmigung des Autors 
in irgendeiner Form reproduziert, gespeichert oder übertragen werden, 
weder elektronisch, mechanisch, durch Fotokopien, Aufnahmen noch auf andere Weise.
\begin{center}\small Printed in Germany\end{center}

\cleardoublepage

\chapter*{Vorwort}
\markboth{Vorwort}{Vorwort} % sorgt für die Kopfzeile, optional
% kein addcontentsline → Vorwort erscheint NICHT im Inhaltsverzeichnis



Die Entstehung dieses Buches war von einer tiefen Faszination für das Licht und seinen fundamentalen Vermittler – das 



\begin{flushright}
	\textit{Dipl.-Ing.(FH) Christian Weilharter} \\
	\vspace{0.5em}
	[Traunstein, 2025]
\end{flushright}


